% Chapter 26: Principal Stratification
% A First Course in Causal Inference - Peng Ding
% Textbook Chapter 26 topic

\chapter{主分层：条件因果效应与中介分析的基础}\label{ch:26}

\section{本章导论}

\textbf{主分层（Principal Stratification）}是因果推断中一个深刻而优雅的概念框架，由 \cite{RosenbaumRubin:1983b} 首次提出。它为处理一类根本性的难题提供了统一语言——当处理的施与取决于潜在结果时，因果效应如何定义和识别？

经典的潜在结果框架 $Y_i(1)$ 和 $Y_i(0)$ 依赖于一个隐含假设：**处理的施与独立于潜在结果**（SUTVA 的一部分）。然而在许多实际场景中，这个假设根本不成立。例如：
\begin{itemize}
  \item 在疫苗接种研究中，最终接种疫苗的往往是那些认为自己风险较高的人（而潜在地，正是这些人可能本来就更健康或更不健康）；
  \item 在工作培训项目中，自愿参加培训的个体可能在能力和动机上与未参加者存在系统性差异；
  \item 在番茄酱实验中，购买番茄酱的家庭与不购买的家庭在购买力、品牌偏好等方面存在差异。
\end{itemize}

这些现象的共同特征是：**处理的施与（compliance）取决于潜在结果**——用更专业的术语说，处理并非独立于潜在结果。这种情况下，标准的因果效应定义（ACE、ATT 等）可能产生令人困惑的解释。

主分层框架的核心思想是：根据\textbf{潜在结果的联合分层}（joint stratification of potential outcomes）对总体进行划分，然后在每一层内讨论因果效应。每一层内部的单元在处理施与上具有相同的"行为模式"，从而使得在该层内的因果推断更加清晰。

\section{主分层的定义与记号}

考虑二元处理 $Z_i \in \{0, 1\}$。对每个单元 $i$，定义两个层面的潜在结果：
\begin{itemize}
  \item \textbf{处理施与层面的潜在结果}：$D_i(1)$ 和 $D_i(0)$，分别表示若强制给予处理或对照时的处理指示；
  \item \textbf{结果层面的潜在结果}：$Y_i(1)$ 和 $Y_i(0)$，分别表示若实际收到处理或对照时的结果。
\end{itemize}

注意：$D_i(1)$ 和 $D_i(0)$ 是关于\textbf{处理施与}的潜在结果，而非关于结果本身。这意味着处理施与本身也可以被视为一个因果变量。

根据 $D_i(1)$ 和 $D_i(0)$ 的取值，我们可以将总体划分为四个\textbf{主层（principal stratum）}：
\[
\mathcal{S}_c = \{i : D_i(1) = D_i(0) = c\}, \quad c \in \{0, 1\},
\]
和
\[
\mathcal{S}_C = \{i : D_i(1) = 1, D_i(0) = 0\},
\]
其中 $\mathcal{S}_C$ 是**始终-compliance 层（always-takers）**，$\mathcal{S}_1$ 是**始终-takers 层**，$\mathcal{S}_0$ 是**始终-conformers 层**或**never-takers 层**，$\mathcal{S}_N = \{i : D_i(1) = 0, D_i(0) = 1\}$ 是**never-takers 层**。

\textbf{关键性质}：主层是根据潜在处理施与定义的，因此是处理实施前就已经决定的内在特征分布。我们永远无法从观测数据中直接识别各层的规模（因为这要求我们知道 $D_i(1)$ 和 $D_i(0)$），但可以通过因果假设（如单调性）来间接推断。

\section{核心因果 estimand：主层平均因果效应}\label{sec:26-place-ace}

主分层框架的核心贡献之一是提出了更精细的因果效应定义。

\textbf{主层平均因果效应（Principal Stratified Average Causal Effect, PSACE）}：
\[
\tau^{\mathcal{S}_s} = \mathbb{E}\bigl( Y_i(1) - Y_i(0) \mid i \in \mathcal{S}_s \bigr),
\]
即在主层 $\mathcal{S}_s$ 内的平均因果效应。

不同的 estimand 有不同的解释：
\begin{itemize}
  \item $\tau^{\mathcal{S}_C}$：对\textbf{依从者层（compliers）}的平均因果效应（Complier Average Causal Effect, CACE）——这正是 IV 框架中识别的效应！
  \item $\tau^{\mathcal{S}_1}$：对始终-takers 的平均因果效应
  \item $\tau^{\mathcal{S}_0}$：对 never-takers 的平均因果效应
\end{itemize}

这揭示了一个深刻结论：**工具变量（IV）方法识别的恰好是 compliers 层上的因果效应**——这就是为什么 IV 的解释依赖于依从性假设。

\section{主分层在中介分析中的应用}\label{sec:26-mediation}

主分层框架在中介分析中也发挥核心作用。考虑一个两步处理设定：
\[
Z \rightarrow M \rightarrow Y,
\]
其中 $Z$ 是处理，$M$ 是中介，$Y$ 是结果。

在这个框架中，我们关注一个微妙的问题：$M$ 的变化是否完全由 $Z$ 引起？如果某些个体的 $M$ 在 $Z=0$ 时已经处于高水平（无论 $Z$ 如何），那么 $Z$ 对 $M$ 的"因果效应"在这些个体上为零。

\textbf{因果中介分析中的主分层}：根据 $M_i(1)$ 和 $M_i(0)$ 对个体进行分层（这与上面的 $D$ 分层框架完全平行），可以得到：
\begin{itemize}
  \item \textbf{直接效应}：$Y_i(1, M_i(0)) - Y_i(0, M_i(0))$，即在反事实 $M_i(0)$ 下，$Z$ 的变化对 $Y$ 的影响；
  \item \textbf{间接效应}：$Y_i(1, M_i(1)) - Y_i(1, M_i(0))$，即通过 $M$ 传导的效应。
\end{itemize}

\section{可识别性与关键假设}\label{sec:26-identify}

主层本身是不可识别的（因为我们无法观测 $D_i(1)$ 和 $D_i(0)$ 同时取值）。但通过额外的因果假设，我们可以识别某些主层的因果效应：

\textbf{单调性假设}：$D_i(1) \geq D_i(0)$ 对所有 $i$ 成立。这意味着没有人会因为处理而反而放弃处理——这是一个弱的非稀释假设。

在单调性下，始终-takers 层（$\mathcal{S}_1$）的规模为：
\[
\mathbb{E}(\mathcal{S}_1) = \mathbb{E}(D_i(1) - D_i(0)) = \mathbb{E}(D_i(1)) - \mathbb{E}(D_i(0)),
\]
即处理组和对照组的处理施与率之差。

\textbf{排他性假设（Exclusion Restriction）}：对 compliers 层，直接处理效应为零——这意味着处理 $Z$ 对结果 $Y$ 的影响完全通过处理指示变量 $D$ 实现。

这些假设共同构成了 IV 识别策略的基础。

\section{本章小结}\label{sec:26-summary}

主分层框架为因果推断提供了一个统一的概念工具，使得我们能够：
\begin{enumerate}
  \item 在处理施与取决于潜在结果的复杂场景中，仍然可以谈论因果效应；
  \item 将 IV 方法的估计量解释为 compliers 层上的因果效应；
  \item 在中介分析中精确区分直接效应和间接效应；
  \item 通过透明的假设（单调性、排他性）明确识别策略的适用范围。
\end{enumerate}

主分层是现代因果推断的一块基石，建议后续阅读 \cite{RosenbaumRubin:1983b} 的原文以深入理解其动机和数学结构。

\parNotes{经典参考文献：\cite{RosenbaumRubin:1983b}（《\textit{中华流行病学杂志}》），以及 \cite{AngristPischke:2009} 第 4 章关于 IV 和主分层的讨论。}
